\section{Appendix}
\label{sec:appendix}

\subsection{Prediction Prompt Templates}
All prediction queries used temperature $= 0.0$, \texttt{max\_tokens} $= 10$, and seed $= 42$. The augmentation line is omitted for C0.

\subsubsection*{Primary Ablation (BBBP, BACE):}
\begin{verbatim}
Given this molecule (SMILES: {smiles})
[Additional information: {augmentation}]

On a scale of 0 to 1, what is the probability
that this molecule can {task_description}?
Respond with ONLY a number between 0 and 1.
Nothing else.
\end{verbatim}
Task descriptions: ``penetrate the blood-brain barrier'' (BBBP), ``inhibit the BACE-1 enzyme'' (BACE).

\subsubsection*{Extended Benchmark: Tox21 (Classification):}
\begin{verbatim}
Given this molecule (SMILES: {smiles})
[Additional information: {augmentation}]

Is this molecule toxic for the NR-AR pathway?
Respond with ONLY a number representing the
probability of toxicity (between 0 and 1).
Nothing else.
\end{verbatim}

\subsubsection*{Extended Benchmark: ESOL (Regression):}
\begin{verbatim}
Given this molecule (SMILES: {smiles})
[Additional information: {augmentation}]

What is the measured log solubility in mols
per litre for this molecule?
Respond with ONLY a numerical value (e.g.,
-2.5). Nothing else.
\end{verbatim}

Hallucination generation used temperature $= 0.9$, \texttt{max\_tokens} $= 150$, with the following type-specific prompts.

\subsubsection*{Primary Ablation Generation Prompts (BBBP, BACE):}

\textbf{C3 (Free Hallucination):}
\textit{``You are a speculative chemist. Given SMILES: \{smiles\}. Generate a creative, scientifically-sounding description of this molecule's potential properties and biological mechanisms. Be creative and speculative. Write 2--3 sentences only.''}

\textbf{C4a (Structural Augmentation):}
\textit{``Given SMILES: \{smiles\}. Describe the structural features, functional groups, and atomic composition. Be specific and speculative.''}

\textbf{C4b (Property Inversion):}
\textit{``Given SMILES: \{smiles\}. Describe its solubility, lipophilicity, and membrane permeability with specific values.''}

\textbf{C4c (Mechanism Fabrication):}
\textit{``Given SMILES: \{smiles\}. Hypothesize biological targets and cellular mechanisms. Be specific about protein targets.''}

\subsubsection*{Extended Benchmark Generation Prompts (Tox21, ESOL):}
\textbf{Note:} The generation prompts for the extended benchmarks were adapted from the primary ablation prompts. While designed to serve the same functional role, the exact phrasing differs, which constitutes a minor methodological confound acknowledged in the main text.

\textbf{C3 (Free Hallucination):}
\textit{``Provide a detailed scientific description of the molecular structure and properties for the molecule with SMILES: \{smiles\}. Focus on its potential biological interactions and physicochemical characteristics.''}

\textbf{C4a (Structural Augmentation):}
\textit{``Describe the scaffold and structural features of the molecule with SMILES: \{smiles\}. Mention specific rings, functional groups, and their spatial arrangement.''}

\subsection{Control Condition Templates}

\subsubsection*{C2 (Chemical Priming):}
Template-generated text using scientific vocabulary:
\textit{``This compound exhibits \{adj\} structural coherence with \{noun\} interaction domains, demonstrating \{adj\} physicochemical characteristics under \{noun\} conditions.''}

Adjective pool: moderate, intermediate, variable, standard, conventional, typical.
Noun pool: hypothetical, theoretical, computational, structural, molecular, chemical.

\textbf{Note:} This template contains chemical vocabulary (``structural'', ``molecular'', ``physicochemical''), making C2 a \emph{partial} stylistic control. See C2b for the stricter control.

\subsubsection*{C2b (Pure Gibberish):}
Template-generated text using exclusively non-chemical vocabulary:
\textit{``This entity exhibits \{adj\} portfolio coherence with \{noun\} governance domains, demonstrating \{adj\} bureaucratic characteristics under \{noun\} conditions.''}

Adjective pool: variable, standard, conventional, typical, moderate, intermediate.
Noun pool: portfolio, transactional, governance, bureaucratic, administrative, regulatory.

Vocabulary audit confirmed 0\% overlap with a curated chemical/biological term list.

\subsection{Taxonomy Classifier Pseudocode}

\begin{enumerate}
    \item \textbf{SP Detection:} For each element in \{F, Cl, Br, I, P, S, N\}, check if associated keywords appear in the generated text. Cross-reference with \texttt{mol.GetAtoms()} to verify element absence. If a keyword matches an absent element, classify as SP.
    \item \textbf{PI Detection:} Compute MolLogP via RDKit. If LogP $> 3$ and text contains ``water soluble'', ``hydrophilic'', or ``highly soluble'', classify as PI. If LogP $< 1$ and text contains ``lipophilic'', ``hydrophobic'', or ``membrane permeable'', classify as PI.
    \item \textbf{MF Detection:} Keyword matching for 9 biological targets: ACE2, HER2, EGFR, VEGF, p53, BCL-2, COX-2, dopamine, serotonin.
    \item \textbf{CC (Default):} If none of the above conditions trigger, classify as Contextual Confabulation.
\end{enumerate}

\subsection{Representative Hallucination Examples}
\begin{itemize}
    \item \textbf{SP Example (SMILES: c1ccccc1, Benzene):} ``Contains a reactive chlorine atom that enhances metabolic stability.'' (Benzene contains no chlorine atoms.)
    \item \textbf{PI Example (SMILES: high-LogP compound):} ``This highly water-soluble molecule readily dissolves in aqueous media.'' (Contradicts computed LogP $> 3$.)
    \item \textbf{MF Example (Diazepam):} ``Acts as a potent inhibitor of the VEGF protein kinase receptor.'' (Diazepam is a GABA receptor modulator, not a VEGF inhibitor.)
    \item \textbf{CC Example:} ``This molecule shows promising drug-like properties and may have therapeutic applications in treating various diseases.'' (Generic, non-falsifiable.)
\end{itemize}

\subsection{API Failure Rates}\label{subsec:api_failures}
Across the BBBP experiment (204 molecules $\times$ 9 conditions $= 1{,}836$ predictions), API failures (rate limits, parsing errors) that defaulted to $p = 0.5$ accounted for less than 2\% of total predictions. The exact count is logged in the experiment checkpoint files.

\subsection{Implementation Environment}
Experiments were executed on a local system (Windows 11) using the Groq API for LLM inference. Total runtime for the pilot study was approximately 4.5 hours. Library versions: \texttt{rdkit-pypi} v2023.9.1~\cite{landrum2023rdkit}, \texttt{deepchem} v2.7.1~\cite{ramsundar2019deepchem}, \texttt{groq} Python SDK, \texttt{scipy} v1.11, \texttt{scikit-learn} v1.3.
